\section*{План работы}

\begin{table}[h!]
\begin{center}
\begin{tabular}{|m{28em} | m{5em}|}
\hline
\textbf{Задача} & \textbf{Дата}\\
\hline
\hline
 Разработать архитектуру MVP системы. Реализовать базовый планировщик задач и логику Telegram-бота с in-memory хранилищем. & 10/11/2025 \\
\hline
 Спроектировать реляционную схему БД. Интегрировать PostgreSQL (Liquibase) и реализовать слой доступа к данным (ORM и SQL). & 30/11/2025 \\
\hline
 Внедрить асинхронное межсервисное взаимодействие (Apache Kafka) и систему кэширования данных (Redis). & 15/12/2025 \\
\hline
 Повысить отказоустойчивость распределенной системы: реализовать паттерны Timeout, Retry, Rate Limiter и Circuit Breaker. & 31/12/2025 \\
\hline
 Настроить инфраструктуру наблюдаемости (Prometheus, Grafana). Внедрить бизнес-метрики. Контейнеризовать микросервисы (Docker). & 15/01/2026 \\
\hline
 Расширить систему сбора данных анонсов алгоритмами парсинга Web-ресурсов и реализовать синхронную валидацию источников. & 31/01/2026 \\
\hline
 Интегрировать YandexGPT API для интеллектуального анализа, структурирования и суммаризации анонсов. & 04/02/2026 \\
\hline
 Спроектировать подсистему идентификации пользователей. Реализовать Yandex OAuth и связывание аккаунтов в едином цифровом профиле пользователя. & 20/02/2026 \\
\hline
 Разработать механизм типизированной фильтрации и механизм подбора соотвествующих фильтрам событий для пользователей. & 28/02/2026 \\
\hline
 Внедрить маршрутизатор омниканальных уведомлений (Notification Router) и подсистему E-mail рассылок на базе HTML-шаблонов. & 10/03/2026 \\
\hline
 Инициализировать Web-клиент (Next.js, App Router). Разработать базовые компоненты интерфейса и флоу JWT-авторизации. & 20/03/2026 \\
\hline
 Разработать интерактивный дашборд и модуль управления подписками на источники с строгой клиентской валидацией (Zod) и Server Actions. & 15/04/2026 \\
\hline
\end{tabular}
\caption{План работ по разработке системы TechEventsHub.}
\label{tab:work_plan}
\end{center}
\end{table}

Итоговым результатом выполнения описанного плана стал готовый программный продукт. 
Исходный код разработанной системы размещен в публичном репозитории на платформе GitHub по ссылке: \url{https://github.com/ivanz851/tech_events_hub_platform}.
Инструкция по локальному развертыванию системы приведена в Приложении~\hiddenref{app:local-app-deployment}.